\documentclass[12pt]{article}
\usepackage[utf8]{inputenc}
\usepackage[T1]{fontenc}
\usepackage{hyperref}
\title{Korubo Language Resources for K-12 Curriculum Development}
\author{Compiled by Claude Sonnet 5.5 (claude-sonnet-5-5) for Andr\'as Kornai}
\date{}
\begin{document}
\maketitle
\section{Introduction}
Korubo (ISO 639-3 xor; Glottocode koru1247) is a Panoan language of the Javari Valley in Amazonas, Brazil. The Korubo were contacted by FUNAI in 1996, and many remain in isolation. The population is listed as 127 (Funai, 2020). The search results showed no teaching materials. Because many Korubo remain in voluntary isolation, outsiders should treat any curriculum work with great care and defer to the people and FUNAI.

\section{Primary Educational Resources}
None found in the search results.

\section{Linguistic and Reference Materials}
\begin{itemize}
\item A survey of the languages of the Javari River Valley, Brazil (SIL). It concludes that Korubo and Kulina Pano are in Northern Panoan. \url{https://www.silbrazil.org/resources/archives/9091}
\item A UFAM dissertation on the Korubo of the Vale do Javari: \url{https://www.tede.ufam.edu.br/handle/tede/6132}
\item Wikipedia, Korubo: \url{https://en.wikipedia.org/wiki/Korubo}
\end{itemize}

\section{Religious and Community Resources}
Povos Ind\'igenas no Brasil (ISA) page on the Korubo: \url{https://pib.socioambiental.org/en/Povo:Korubo}. The Joshua Project page gives people-group data only: \url{https://joshuaproject.net/people_groups/12805}.

\section{Digital and Audio Resources}
None found in the search results.

\section{Curriculum Development Context}
No known school programme or published orthography was found. Any curriculum would need community consent and have to start from very limited documentation.
\end{document}